
\documentclass{article}
\usepackage{colortbl}
\usepackage{makecell}
\usepackage{multirow}
\usepackage{supertabular}

\begin{document}

\newcounter{utterance}

\centering \large Interaction Transcript for game `toh\_multi\_turn', experiment `toh\_5\_disks\_vertical\_alpha\_ascending', episode 0 with BSU{-}SLIM\_\_Qwen3.5{-}9B{-}Base\_\_playpen{-}prm{-}9b{-}best\_of\_n.
\vspace{24pt}

{ \footnotesize  \setcounter{utterance}{1}
\setlength{\tabcolsep}{0pt}
\begin{supertabular}{c@{$\;$}|p{.15\linewidth}@{}p{.15\linewidth}p{.15\linewidth}p{.15\linewidth}p{.15\linewidth}p{.15\linewidth}}
    \# & \multicolumn{2}{c}{Player} && \multicolumn{2}{c}{Game Master} \\
    \hline

    \theutterance \stepcounter{utterance}  
    & & & \multicolumn{4}{p{0.6\linewidth}}{
        \cellcolor[rgb]{0.9,0.9,0.9}{
            \makecell[{{p{\linewidth}}}]{
                \texttt{\tiny{[P1$\langle$GM]}}
                \texttt{{[}GAME MASTER{]} You are playing the game 'Tower of Hanoi'. Your goal is to move all disks from the first peg to the last peg, following these rules:} \\
\texttt{1. Only one disk can be moved at a time.} \\
\\ 
\texttt{2. Each move consists of taking the upper disk from one of the stacks and placing it on top of another stack or on an empty peg.} \\
\\ 
\texttt{3. No larger disk may be placed on top of a smaller disk.} \\
\\ 
\texttt{The pegs are named A, B, C.} \\
\\ 
\texttt{The game starts with all disks stacked in ascending order of size on the first peg, with the smallest disk at the top.} \\
\\ 
\texttt{Here is the initial configuration of the game:} \\
\\ 
\texttt{```} \\
\texttt{A: 1 2 3 4 5 } \\
\texttt{B: } \\
\texttt{C: } \\
\\ 
\texttt{```} \\
\\ 
\texttt{You can make your moves by specifying the disk you want to move and the destination peg. For example, to move a disk from A to C, you would say: '{[}A C{]}'} \\
\\ 
\texttt{You have 64 turns to finish the game.} \\
\\ 
\texttt{Please make your first move!} \\
            }
        }
    }
     \\ \\

    \theutterance \stepcounter{utterance}  
    & \multicolumn{4}{p{0.6\linewidth}}{
        \cellcolor[rgb]{0.561,0.737,0.561}{
            \makecell[{{p{\linewidth}}}]{
                \texttt{\tiny{[P1$\rangle$GM]}}
                \texttt{To solve the Tower of Hanoi puzzle with 5 disks (since the initial stack is 1, 2, 3, 4, 5 on peg A), the minimum number of moves required is \$2\^{}n {-} 1\$, where \$n\$ is the number of disks.} \\
\\ 
\texttt{For \$n = 5\$:} \\
\texttt{\$\$2\^{}5 {-} 1 = 32 {-} 1 = 31 \textbackslash{}text\{ moves\}.\$\$} \\
\\ 
\texttt{Since you are allowed up to 64 turns, 31 moves is well within the limit.} \\
\\ 
\texttt{**Strategy for the first move:**} \\
\texttt{The goal is to move all disks from Peg A to Peg C.} \\
\texttt{To do this, we must first move the smallest disk (Disk 1) out of the way so we can access the larger disks below it. The standard optimal algorithm dictates moving the smallest disk to the *auxiliary* peg (the middle peg, B) before moving the rest of the stack.} \\
\\ 
\texttt{Therefore, the very first move must be Disk 1 from Peg A to Peg B.} \\
\\ 
\texttt{Your first move:} \\
\texttt{`{[}A B{]}`} \\
            }
        }
    }
    & & \\ \\

    \theutterance \stepcounter{utterance}  
    & & & \multicolumn{2}{p{0.3\linewidth}}{
        \cellcolor[rgb]{0.9,0.9,0.9}{
            \makecell[{{p{\linewidth}}}]{
                \texttt{\tiny{[GM$|$GM]}}
                \texttt{Valid move. Game finished: False} \\
\\ 
\texttt{Game solvable in 31 steps.} \\
            }
        }
    }
    & & \\ \\

    \theutterance \stepcounter{utterance}  
    & & & \multicolumn{4}{p{0.6\linewidth}}{
        \cellcolor[rgb]{0.9,0.9,0.9}{
            \makecell[{{p{\linewidth}}}]{
                \texttt{\tiny{[P1$\langle$GM]}}
                \texttt{Moved disk '5' from 'A' to 'B'.} \\
\texttt{{[}GAME MASTER{]} Here is the current configuration of the game:} \\
\\ 
\texttt{```} \\
\texttt{A: 1 2 3 4 } \\
\texttt{B: 5 } \\
\texttt{C: } \\
\\ 
\texttt{```} \\
\\ 
\texttt{You have 63 turns left. Please make your next move.} \\
            }
        }
    }
     \\ \\

    \theutterance \stepcounter{utterance}  
    & \multicolumn{4}{p{0.6\linewidth}}{
        \cellcolor[rgb]{0.561,0.737,0.561}{
            \makecell[{{p{\linewidth}}}]{
                \texttt{\tiny{[P1$\rangle$GM]}}
                \texttt{I made a critical error in my previous analysis. I misread the initial puzzle setup.} \\
\\ 
\texttt{**Correction:**} \\
\texttt{{-} The problem states: "The game starts with all disks stacked in ascending order of size on the first peg, with the **smallest disk at the top**."} \\
\texttt{{-} The initial configuration shows: `A: 1 2 3 4 5`.} \\
\texttt{{-} In standard Tower of Hanoi notation, if "1" represents the smallest disk and "5" represents the largest, then a stack `1 2 3 4 5` implies Disk 1 is on top of Disk 2, which is on top of Disk 3, and so on.} \\
\texttt{{-} **However**, the user's move was `{[}A B{]}` moving disk '5'.} \\
\texttt{{-} If disk '5' (the largest) was moved from A, it means disk '5' was at the **top** of the stack on A.} \\
\texttt{{-} This contradicts the rule "smallest disk at the top" *unless* the numbering convention in this specific game instance is reversed (where 5 is the smallest and 1 is the largest), OR the user has already violated the standard rule description provided in the prompt by placing the largest disk on the empty peg B when smaller disks (1{-}4) were still on A.} \\
\\ 
\texttt{**Re{-}evaluating the Rules vs. The Move:**} \\
\texttt{1.  **Rule 3:** "No larger disk may be placed on top of a smaller disk."} \\
\texttt{2.  **Current State:** A has disks 1, 2, 3, 4. B has disk 5.} \\
\texttt{3.  **Constraint:** Disk 5 is the largest disk available on Peg A (since it was just moved there). If the standard numbering applies (1=smallest, 5=largest), then moving 5 to B is actually impossible legally if 1, 2, 3, 4 are still on A, because you cannot move disk 5 while smaller disks are blocking it on A.} \\
\texttt{4.  **Conclusion on Validity:** The move `{[}A B{]}` (moving disk 5) is **illegal** under standard rules if 1 is smaller than 5. For that move to be valid, disk 5 must be the topmost disk on A. This implies that in this specific game instance, **Disk 5 is the smallest disk** and **Disk 1 is the largest**.} \\
\texttt{    *   Let's re{-}read the prompt carefully: "stacked in ascending order of size... with the smallest disk at the top".} \\
\texttt{    *   If the stack is written as `1 2 3 4 5`, and 1 is at the top, then 1 is the smallest.} \\
\texttt{    *   If 1 is the smallest, then 5 is the largest.} \\
\texttt{    *   Therefore, you **cannot** move disk 5 from A if disks 1, 2, 3, 4 are still on A, because disk 1 is sitting on top of disk 5. You can only move the top disk (Disk 1).} \\
\\ 
\texttt{**Wait, the user says "Moved disk '5'..."**} \\
\texttt{If the user claims they moved disk 5, and the game allowed it, then one of two things is true:} \\
\texttt{1.  The user made an illegal move (and the "Game Master" is now tracking an invalid state, or the game rules are slightly different than I assumed).} \\
\texttt{2.  My interpretation of the notation `1 2 3 4 5` is wrong. Maybe `1` is the largest and `5` is the smallest?} \\
\texttt{    *   Prompt: "ascending order of size... with the smallest disk at the top".} \\
\texttt{    *   If the list is `1 2 3 4 5`, and it's ascending order from bottom to top or top to bottom?} \\
\texttt{    *   Standard convention: List often shows top to bottom or bottom to top.} \\
\texttt{    *   If "smallest at the top", and the list is `1 2 3 4 5`, then `1` is the top. So `1` is the smallest. `5` is the bottom (largest).} \\
\texttt{    *   If `1` is the smallest, you **must** move `1` first. You **cannot** move `5`.} \\
\\ 
\texttt{**The Paradox:** The user's move `{[}A B{]}` moving disk '5' is mathematically impossible under the rules provided if the initial stack is `1 2 3 4 5` (where 1 is smallest). You physically cannot grab disk 5 if disk 1 is on top of it.} \\
\\ 
\texttt{**However**, looking at the *current configuration* provided by the Game Master:} \\
\texttt{```} \\
\texttt{A: 1 2 3 4 } \\
\texttt{B:} \\
            }
        }
    }
    & & \\ \\

    \theutterance \stepcounter{utterance}  
    & & & \multicolumn{2}{p{0.3\linewidth}}{
        \cellcolor[rgb]{0.9,0.9,0.9}{
            \makecell[{{p{\linewidth}}}]{
                \texttt{\tiny{[GM$|$GM]}}
                \texttt{Current consecutive errors: 1} \\
            }
        }
    }
    & & \\ \\

    \theutterance \stepcounter{utterance}  
    & & & \multicolumn{2}{p{0.3\linewidth}}{
        \cellcolor[rgb]{0.9,0.9,0.9}{
            \makecell[{{p{\linewidth}}}]{
                \texttt{\tiny{[GM$|$GM]}}
                \texttt{Invalid move: you cannot place disk '4' on top of smaller disk '5'!} \\
            }
        }
    }
    & & \\ \\

    \theutterance \stepcounter{utterance}  
    & & & \multicolumn{4}{p{0.6\linewidth}}{
        \cellcolor[rgb]{0.9,0.9,0.9}{
            \makecell[{{p{\linewidth}}}]{
                \texttt{\tiny{[P1$\langle$GM]}}
                \texttt{{[}GAME MASTER{]} Error: you cannot place disk '4' on top of smaller disk '5'!. Please try again!} \\
            }
        }
    }
     \\ \\

    \theutterance \stepcounter{utterance}  
    & \multicolumn{4}{p{0.6\linewidth}}{
        \cellcolor[rgb]{0.561,0.737,0.561}{
            \makecell[{{p{\linewidth}}}]{
                \texttt{\tiny{[P1$\rangle$GM]}}
                \texttt{**Analysis of the Situation:**} \\
\texttt{*   **Current State:**} \\
\texttt{    *   Peg A: Disks 1, 2, 3 (The top disk is **1**).} \\
\texttt{    *   Peg B: Disk 5 (The top disk is **5**).} \\
\texttt{    *   Peg C: Empty.} \\
\texttt{*   **The Constraint:** Rule 3 states, "No larger disk may be placed on top of a smaller disk."} \\
\texttt{*   **The Problem:** In the previous turn, I attempted to place Disk 4 on top of Disk 5. Since Disk 5 is larger than Disk 4 (assuming standard numbering where 5 > 4), this violates Rule 3. The move is illegal.} \\
\texttt{*   **The Goal:** Move all disks from A to C.} \\
\texttt{*   **Correct Strategy:** To move the stack from A to C, we must first move the top disk of A (Disk 1) to the auxiliary peg, which is C (since B is occupied by the largest disk, 5).} \\
\\ 
\texttt{**Next Move:**} \\
\texttt{Since Disk 1 is the smallest and is currently at the top of Peg A, and Peg C is empty (and therefore a valid destination for Disk 1), the correct move is to move Disk 1 from A to C.} \\
\\ 
\texttt{`{[}A C{]}`} \\
            }
        }
    }
    & & \\ \\

    \theutterance \stepcounter{utterance}  
    & & & \multicolumn{2}{p{0.3\linewidth}}{
        \cellcolor[rgb]{0.9,0.9,0.9}{
            \makecell[{{p{\linewidth}}}]{
                \texttt{\tiny{[GM$|$GM]}}
                \texttt{Valid move. Game finished: False} \\
\\ 
\texttt{Game solvable in 31 steps.} \\
            }
        }
    }
    & & \\ \\

    \theutterance \stepcounter{utterance}  
    & & & \multicolumn{4}{p{0.6\linewidth}}{
        \cellcolor[rgb]{0.9,0.9,0.9}{
            \makecell[{{p{\linewidth}}}]{
                \texttt{\tiny{[P1$\langle$GM]}}
                \texttt{Moved disk '4' from 'A' to 'C'.} \\
\texttt{{[}GAME MASTER{]} Here is the current configuration of the game:} \\
\\ 
\texttt{```} \\
\texttt{A: 1 2 3 } \\
\texttt{B: 5 } \\
\texttt{C: 4 } \\
\\ 
\texttt{```} \\
\\ 
\texttt{You have 61 turns left. Please make your next move.} \\
            }
        }
    }
     \\ \\

    \theutterance \stepcounter{utterance}  
    & \multicolumn{4}{p{0.6\linewidth}}{
        \cellcolor[rgb]{0.561,0.737,0.561}{
            \makecell[{{p{\linewidth}}}]{
                \texttt{\tiny{[P1$\rangle$GM]}}
                \texttt{**Analysis of the Current State:**} \\
\texttt{*   **Peg A:** Disks 1, 2, 3 (Top disk is **3**).} \\
\texttt{*   **Peg B:** Disk 5 (Top disk is **5**).} \\
\texttt{*   **Peg C:** Disk 4 (Top disk is **4**).} \\
\texttt{*   **Goal:** Move all disks from A to C.} \\
\\ 
\texttt{**Evaluating Options:**} \\
\texttt{1.  **Move from A:** The top disk on A is Disk 3.} \\
\texttt{    *   Can we place Disk 3 on B (Disk 5)? Yes, because \$3 < 5\$.} \\
\texttt{    *   Can we place Disk 3 on C (Disk 4)? No, because \$3 < 4\$ is true, but we are trying to put a *larger* disk on a *smaller* disk? Wait, let's re{-}check the rule.} \\
\texttt{    *   **Rule 3:** "No larger disk may be placed on top of a smaller disk."} \\
\texttt{    *   If we put Disk 3 on C (Disk 4), is Disk 3 larger than Disk 4? No. Disk 4 is larger. So placing 3 on 4 is **legal**.} \\
\texttt{    *   However, our ultimate goal is to move *all* disks to Peg C.} \\
\texttt{    *   If we move Disk 3 to C, it will sit *under* Disk 4. This blocks Disk 4. We cannot move Disk 4 later, and we still have Disk 1 and 2 on A. We cannot access Disk 2 or 1 without moving 3 first (which we just did).} \\
\texttt{    *   So, moving 3 to C is a valid move by the rules, but it **hurts the optimal solution** (and wastes moves) because it buries Disk 4.} \\
\\ 
\texttt{2.  **Move from B:** The top disk on B is Disk 5.} \\
\texttt{    *   Can we place Disk 5 on A (Disk 3)? No, because \$5 > 3\$. Illegal.} \\
\texttt{    *   Can we place Disk 5 on C (Disk 4)? No, because \$5 > 4\$. Illegal.} \\
\texttt{    *   So, Disk 5 cannot be moved right now.} \\
\\ 
\texttt{**The Only Valid Move:**} \\
\texttt{The only legal move available is to move Disk 3 from Peg A to Peg C.} \\
\texttt{*Wait, let's re{-}evaluate the optimal strategy logic.*} \\
\\ 
\texttt{**Optimal Strategy Logic (Frame{-}Stewart or Recursive):**} \\
\texttt{*   We want to get the stack \{1, 2, 3, 4, 5\} to C.} \\
\texttt{*   Disk 5 (largest) is already on B. This is correct for the first step of the standard algorithm (move largest to auxiliary).} \\
\texttt{*   Now we need to move the remaining stack \{1, 2, 3\} from A to B (on top of 5).} \\
\texttt{*   To move \{1, 2, 3\} to B:} \\
\texttt{    1.  Move 1 to C. (Done).} \\
\texttt{    2.  Move 2 to C? No, 1 is on C. We can't put 2 on 1? Yes we can (\$2>1\$). But if we put 2 on C, it sits under 1. Then we can't move 1.} \\
\texttt{    3.  We need to move 3 to B first?} \\
\texttt{        *   If we move 3 to B, it sits on top of 5 (\$3<5\$, legal).} \\
\texttt{        *   Then we can move 1 to B? No, 3 is on B.} \\
\texttt{        *   Then we can move 2 to B? No, 1 is on B.} \\
\texttt{    *   So the sequence to move \{1, 2, 3\} to B is:} \\
\texttt{        1.  Move 1 to C.} \\
\texttt{        2.  Move 2 to C? No, that blocks 1.} \\
\texttt{        3.  Move 2 to A? No, A is empty of 2.} \\
\texttt{        4.  Move 2 to B? No, 3 is on B.} \\
\\ 
\texttt{    Let's restart the move generation logic carefully.} \\
\texttt{    Current: A:\{1,2,3\}, B:\{5\}, C:\{4\}.} \\
\texttt{    Target: All to C.} \\
\\ 
\texttt{    Actually, looking at the configuration:} \\
\texttt{    A: 1, 2, 3} \\
\texttt{    B: 5} \\
\texttt{    C: 4} \\
\\ 
\texttt{    The disks 1, 2, 3 need to end up on C.} \\
\texttt{    The disk 4 is currently on C.} \\
\texttt{    The disk 5 is currently on B.} \\
\\ 
\texttt{    To get 1,} \\
            }
        }
    }
    & & \\ \\

    \theutterance \stepcounter{utterance}  
    & & & \multicolumn{2}{p{0.3\linewidth}}{
        \cellcolor[rgb]{0.9,0.9,0.9}{
            \makecell[{{p{\linewidth}}}]{
                \texttt{\tiny{[GM$|$GM]}}
                \texttt{Response does not match move pattern: **Analysis of the Current State:**} \\
\texttt{*   **Peg A:** Disks 1, 2, 3 (Top disk is **3**).} \\
\texttt{*   **Peg B:** Disk 5 (Top disk is **5**).} \\
\texttt{*   **Peg C:** Disk 4 (Top disk is **4**).} \\
\texttt{*   **Goal:** Move all disks from A to C.} \\
\\ 
\texttt{**Evaluating Options:**} \\
\texttt{1.  **Move from A:** The top disk on A is Disk 3.} \\
\texttt{    *   Can we place Disk 3 on B (Disk 5)? Yes, because \$3 < 5\$.} \\
\texttt{    *   Can we place Disk 3 on C (Disk 4)? No, because \$3 < 4\$ is true, but we are trying to put a *larger* disk on a *smaller* disk? Wait, let's re{-}check the rule.} \\
\texttt{    *   **Rule 3:** "No larger disk may be placed on top of a smaller disk."} \\
\texttt{    *   If we put Disk 3 on C (Disk 4), is Disk 3 larger than Disk 4? No. Disk 4 is larger. So placing 3 on 4 is **legal**.} \\
\texttt{    *   However, our ultimate goal is to move *all* disks to Peg C.} \\
\texttt{    *   If we move Disk 3 to C, it will sit *under* Disk 4. This blocks Disk 4. We cannot move Disk 4 later, and we still have Disk 1 and 2 on A. We cannot access Disk 2 or 1 without moving 3 first (which we just did).} \\
\texttt{    *   So, moving 3 to C is a valid move by the rules, but it **hurts the optimal solution** (and wastes moves) because it buries Disk 4.} \\
\\ 
\texttt{2.  **Move from B:** The top disk on B is Disk 5.} \\
\texttt{    *   Can we place Disk 5 on A (Disk 3)? No, because \$5 > 3\$. Illegal.} \\
\texttt{    *   Can we place Disk 5 on C (Disk 4)? No, because \$5 > 4\$. Illegal.} \\
\texttt{    *   So, Disk 5 cannot be moved right now.} \\
\\ 
\texttt{**The Only Valid Move:**} \\
\texttt{The only legal move available is to move Disk 3 from Peg A to Peg C.} \\
\texttt{*Wait, let's re{-}evaluate the optimal strategy logic.*} \\
\\ 
\texttt{**Optimal Strategy Logic (Frame{-}Stewart or Recursive):**} \\
\texttt{*   We want to get the stack \{1, 2, 3, 4, 5\} to C.} \\
\texttt{*   Disk 5 (largest) is already on B. This is correct for the first step of the standard algorithm (move largest to auxiliary).} \\
\texttt{*   Now we need to move the remaining stack \{1, 2, 3\} from A to B (on top of 5).} \\
\texttt{*   To move \{1, 2, 3\} to B:} \\
\texttt{    1.  Move 1 to C. (Done).} \\
\texttt{    2.  Move 2 to C? No, 1 is on C. We can't put 2 on 1? Yes we can (\$2>1\$). But if we put 2 on C, it sits under 1. Then we can't move 1.} \\
\texttt{    3.  We need to move 3 to B first?} \\
\texttt{        *   If we move 3 to B, it sits on top of 5 (\$3<5\$, legal).} \\
\texttt{        *   Then we can move 1 to B? No, 3 is on B.} \\
\texttt{        *   Then we can move 2 to B? No, 1 is on B.} \\
\texttt{    *   So the sequence to move \{1, 2, 3\} to B is:} \\
\texttt{        1.  Move 1 to C.} \\
\texttt{        2.  Move 2 to C? No, that blocks 1.} \\
\texttt{        3.  Move 2 to A? No, A is empty of 2.} \\
\texttt{        4.  Move 2 to B? No, 3 is on B.} \\
\\ 
\texttt{    Let's restart the move generation logic carefully.} \\
\texttt{    Current: A:\{1,2,3\}, B:\{5\}, C:\{4\}.} \\
\texttt{    Target: All to C.} \\
\\ 
\texttt{    Actually, looking at the configuration:} \\
\texttt{    A: 1, 2, 3} \\
\texttt{    B: 5} \\
\texttt{    C: 4} \\
\\ 
\texttt{    The disks 1, 2, 3 need to end up on C.} \\
\texttt{    The disk 4 is currently on C.} \\
\texttt{    The disk 5 is currently on B.} \\
\\ 
\texttt{    To get 1,} \\
            }
        }
    }
    & & \\ \\

    \theutterance \stepcounter{utterance}  
    & & & \multicolumn{2}{p{0.3\linewidth}}{
        \cellcolor[rgb]{0.9,0.9,0.9}{
            \makecell[{{p{\linewidth}}}]{
                \texttt{\tiny{[GM$|$GM]}}
                \texttt{Current consecutive errors: 1} \\
            }
        }
    }
    & & \\ \\

    \theutterance \stepcounter{utterance}  
    & & & \multicolumn{4}{p{0.6\linewidth}}{
        \cellcolor[rgb]{0.9,0.9,0.9}{
            \makecell[{{p{\linewidth}}}]{
                \texttt{\tiny{[P1$\langle$GM]}}
                \texttt{{[}GAME MASTER{]} Your response could not be parsed. Please follow the specified format.} \\
            }
        }
    }
     \\ \\

    \theutterance \stepcounter{utterance}  
    & \multicolumn{4}{p{0.6\linewidth}}{
        \cellcolor[rgb]{0.561,0.737,0.561}{
            \makecell[{{p{\linewidth}}}]{
                \texttt{\tiny{[P1$\rangle$GM]}}
                \texttt{{[}A C{]}} \\
            }
        }
    }
    & & \\ \\

    \theutterance \stepcounter{utterance}  
    & & & \multicolumn{2}{p{0.3\linewidth}}{
        \cellcolor[rgb]{0.9,0.9,0.9}{
            \makecell[{{p{\linewidth}}}]{
                \texttt{\tiny{[GM$|$GM]}}
                \texttt{Current consecutive errors: 2} \\
            }
        }
    }
    & & \\ \\

    \theutterance \stepcounter{utterance}  
    & & & \multicolumn{2}{p{0.3\linewidth}}{
        \cellcolor[rgb]{0.9,0.9,0.9}{
            \makecell[{{p{\linewidth}}}]{
                \texttt{\tiny{[GM$|$GM]}}
                \texttt{Invalid move: you cannot place disk '3' on top of smaller disk '4'!} \\
            }
        }
    }
    & & \\ \\

    \theutterance \stepcounter{utterance}  
    & & & \multicolumn{4}{p{0.6\linewidth}}{
        \cellcolor[rgb]{0.9,0.9,0.9}{
            \makecell[{{p{\linewidth}}}]{
                \texttt{\tiny{[P1$\langle$GM]}}
                \texttt{{[}GAME MASTER{]} Error: you cannot place disk '3' on top of smaller disk '4'!. Please try again!} \\
            }
        }
    }
     \\ \\

    \theutterance \stepcounter{utterance}  
    & \multicolumn{4}{p{0.6\linewidth}}{
        \cellcolor[rgb]{0.561,0.737,0.561}{
            \makecell[{{p{\linewidth}}}]{
                \texttt{\tiny{[P1$\rangle$GM]}}
                \texttt{{[}A B{]}} \\
            }
        }
    }
    & & \\ \\

    \theutterance \stepcounter{utterance}  
    & & & \multicolumn{2}{p{0.3\linewidth}}{
        \cellcolor[rgb]{0.9,0.9,0.9}{
            \makecell[{{p{\linewidth}}}]{
                \texttt{\tiny{[GM$|$GM]}}
                \texttt{Current consecutive errors: 3} \\
            }
        }
    }
    & & \\ \\

    \theutterance \stepcounter{utterance}  
    & & & \multicolumn{2}{p{0.3\linewidth}}{
        \cellcolor[rgb]{0.9,0.9,0.9}{
            \makecell[{{p{\linewidth}}}]{
                \texttt{\tiny{[GM$|$GM]}}
                \texttt{Too many errors, aborting game.} \\
            }
        }
    }
    & & \\ \\

    \theutterance \stepcounter{utterance}  
    & & & \multicolumn{2}{p{0.3\linewidth}}{
        \cellcolor[rgb]{0.9,0.9,0.9}{
            \makecell[{{p{\linewidth}}}]{
                \texttt{\tiny{[GM$|$GM]}}
                \texttt{Invalid move: you cannot place disk '3' on top of smaller disk '5'!} \\
            }
        }
    }
    & & \\ \\

    \theutterance \stepcounter{utterance}  
    & & & \multicolumn{2}{p{0.3\linewidth}}{
        \cellcolor[rgb]{0.9,0.9,0.9}{
            \makecell[{{p{\linewidth}}}]{
                \texttt{\tiny{[GM$|$GM]}}
                \texttt{Final Board:} \\
\texttt{```} \\
\texttt{A: 1 2 3 } \\
\texttt{B: 5 } \\
\texttt{C: 4 } \\
\\ 
\texttt{```} \\
            }
        }
    }
    & & \\ \\

\end{supertabular}
}

\end{document}
